\scanpage{100}
\mspage{\handwritten{$4'$}}
yesterday.

What may look less conviencing though at first sight, is my choice to view \inserted{\handwritten{as non-primitive}} even the ``degeneration maps'' $k_1^i:F_i\to F_{i+1}$, associating to every $i$-object the $(i+1)$-object acting as an identity on the former, In all cases I have met so far, these maps are either given beforehand with the structure (of a $1$-category or $2$-category say), or they can be uniquely deduced from the axioms. In the present set-up however, they seem to me to appear more naturally as ``primary'' (not as primitive) structure, much in the same way as the $\mathrm{inv}_i$. Different choices for associating an $\infty$-groupoid to a topological space, while yielding the same base-sets $F_i$, will however (according to this point of view) give rise to different maps $k_1^i$. The main motivation for this point of view comes from the fact that the mechanism for a uniform construction of the chain of ever higher order structures makes a basic use of the source and target maps only and of their ``transposés'' (see below), and (it seems to me) \unsure{not|act} at all of the degeneration maps, which in this respect rather \handwritten{are} confusing \deleted{a}\inserted{the real} picture\handwritten{,} \inserted{\handwritten{if viewed as ``primitive''.}} - The degeneration maps rather appear as typical cases of primary structure, probably of especial significance in the practical handling of $\infty$-groupoids, but not at all in the conceptual machinery leading up to the con\deleted{\unsure{stru}}struction of the structure species \deleted{of} \handwritten{``}$\infty$-groupoids.\handwritten{''}

- Much in the same way, the composition operations $\overset{i}{*}_1$ are viewed as primary, \deleted{not} as primitive or skeletal structure. Their description for the fundamental $\infty$-groupoid of a space - for instance the description of composition of paths - depends on arbitray choices, such as the choice of an isomorphism (say) between $(I,1)\handwritten{\amalg}_e(I,0)$ and $I$, where $I$ is the unit interval, much in the same way as the notion of an inverse of a path depends on the choice of an isomorphism of $I$ with itself, exchanging the two end-points $0$ and $1$. The operations $\overset{i}{*}_2$ of Godement take sense only once the composition operations $*_1$ are defined - they are ``secondary structure'', and successively the operations $*_3\ldots *_i$ appear as ternary etc structure. This is correctly suggested by the notations which I chose\deleted{\unsure{n}} yesterday, where however \handwritten{I} hastily threw\deleted{\unsure{gh|ht}} all those operations into a same pot baptized ``primary structure''!

It is about time though to come to a tentative precise definition o\unsure{r|f} description of the process of stepwise introduction of an increasing chain of higher order structure. This will be done by introducing a canonical sequence of categories and functors
\[
\underline{C}_0\to\underline{C}_1\to\underline{C}_2\to\ldots\to\underline{C}_n\to\underline{C}_{n+1}\to\unsure{\ldots|-}\ ,
\]
% archive sheet number: 99
where $\underline{C}_n$ denotes the \deleted{``universal''} category harbouring the ``universal'' parti%
